\documentclass[12pt]{article}
\usepackage[utf8]{inputenc}
\usepackage[T1]{fontenc}
\usepackage{hyperref}
\title{Thuri Language Resources for K-12 Curriculum Development}
\author{Compiled by Claude Sonnet 5 (claude-sonnet-5) for Andr\'as Kornai}
\date{}
\begin{document}
\maketitle

\section{Introduction}
Thuri (ISO 639-3 thu) is a language of South Sudan. The search found almost nothing about it: no dictionary, grammar, primer or school materials. The only figure found is 6,600 speakers, reported in 1956.

\section{Primary Educational Resources}
None found.

\section{Linguistic and Reference Materials}
\begin{itemize}
\item Encyclopedic overviews only: \url{https://en.wikipedia.org/wiki/Thuri_language} and \url{https://en.wikipedia.org/wiki/Thuri_people}
\item Wikitongues language page, which lists no resources beyond basic information: \url{https://wikitongues.org/languages/thu/}
\end{itemize}

\section{Religious and Community Resources}
None found.

\section{Digital and Audio Resources}
None found.

\section{Curriculum Development Context}
Thuri should not be confused with Beli (blm), a Central Sudanic language. Nothing usable for a curriculum was found, so any programme would need fieldwork, starting with a check on how many speakers there are today.

\end{document}